\documentclass[journal,twocolumn]{IEEEtran}

% ─── Packages ────────────────────────────────────────────────────────────────
\usepackage[T1]{fontenc}
\usepackage[utf8]{inputenc}
\usepackage[english]{babel}
\usepackage{amsmath}
\usepackage{amssymb}
\usepackage{graphicx}
\usepackage{booktabs}
\usepackage{array}
\usepackage{tabularx}
\usepackage{multirow}
\usepackage{enumitem}
\usepackage{hyperref}
\usepackage{xcolor}
\usepackage{soul}        
\usepackage{microtype}
\usepackage{cite}

\hypersetup{
  colorlinks=true,
  linkcolor=blue,
  citecolor=blue,
  urlcolor=blue
}

% Tighter list spacing to match IEEE style
\setlist[itemize]{noitemsep, topsep=2pt, parsep=0pt, partopsep=0pt}
\setlist[enumerate]{noitemsep, topsep=2pt, parsep=0pt, partopsep=0pt}

% ─── Title & Authors ─────────────────────────────────────────────────────────
\begin{document}

\title{Robotic System Implementation Report for \\ AUTONOMOUS SMART CART}

\author{
  Ashwin Murali Thanalapati (M01037932),~
  Mohammed Shalaby (M01035318),~
  Jayashanka Anushan (M01037028),~
  Vignesh Lakshmanasamy (M01026685)\\
}

\maketitle

% ─── Abstract ────────────────────────────────────────────────────────────────
\begin{abstract}
This paper presents the design, implementation, and evaluation of an
Autonomous Smart Cart system developed using ROS2 and Gazebo simulation.
The system integrates a 2D LiDAR sensor with a follow-me behavior and
obstacle avoidance controller, enabling a mobile robot platform to
autonomously track a user within a simulated indoor retail environment.
The ROS2 Navigation2 (Nav2) stack is used for path planning and
localization, while Gazebo provides a physics-accurate simulation
environment for testing and validation. Initial simulation results in Gazebo demonstrate successful follow-me 
behavior and obstacle avoidance within a simulated indoor retail 
environment. Full quantitative evaluation is ongoing; preliminary 
findings and system architecture are presented in this draft report.
\end{abstract}

\IEEEpeerreviewmaketitle

% ─── I. Introduction ─────────────────────────────────────────────────────────
\section{Introduction}

Shopping malls and large retail environments present significant
navigation challenges for users and robotic systems alike.
Conventional shopping carts require continuous physical effort and
lack intelligent features such as obstacle detection and adaptive
movement, making them inefficient and potentially unsafe,
particularly for elderly users and parents.

This paper presents the implementation of an Autonomous Smart Cart
robotic system, building on the requirements defined in CW1.
The focus of CW2 is on realizing key functionalities including
follow-me behavior, obstacle avoidance, and sensor integration
using ROS2 and Gazebo simulation. The physical hardware prototype
has been replaced by a verified Gazebo simulation environment,
with RViz used for real-time sensor and navigation visualization.

The system integrates sensing, control, and decision-making to enable safe and efficient navigation in retail environments.

The remainder of this paper is structured as follows:
Section~II reviews related work; Section~III describes the system
design and simulation implementation; Section~IV outlines
non-functional requirements; Section~V presents preliminary
experimental results; Section~VI discusses feasibility; and
Section~VII presents the testing specifications and risk assessment;
Section~VIII discusses future improvements; and the appendices cover
the use of generative AI and the design thinking process.


% ─── II. Related Works ───────────────────────────────────────────────────────
\section{Related Works}

Shopping malls and large retail environments are designed to accommodate
high volumes of customers, often resulting in crowded aisles, dynamic
pedestrian flow, and limited maneuvering space. In such environments,
conventional shopping carts present operational challenges including
limited maneuverability, reduced braking control under heavy loads, and
increased physical strain on users. These challenges are particularly
significant for elderly shoppers and parents managing children, where
divided attention and physical effort increase the risk of minor
collisions and unsafe interactions \cite{goodrich2007, ting2003}.

Studies in retail ergonomics show that pushing heavily loaded carts
significantly increases user exertion and reduces mobility efficiency,
particularly among elderly individuals \cite{wirawan2024}. Additionally,
pedestrian flow research shows that dense crowd conditions require
constant micro-adjustments in direction and speed to avoid collisions,
increasing cognitive load and reaction time demands \cite{molnar1995, ishii2015}.
These conditions can lead to sudden stops, braking instability, and
unintended contact with other shoppers.

From a robotics perspective, indoor environments such as shopping malls
present complex navigation challenges. Unlike structured industrial
settings, retail environments are highly dynamic, with unpredictable
human motion, temporary obstacles, and occlusions \cite{thrun2005}.
Human--robot interaction (HRI) research emphasizes the importance of
adaptive speed control, obstacle prediction, and intuitive user feedback
mechanisms to ensure safety and social acceptance in shared environments
\cite{goodrich2007}.

Furthermore, collision-related incidents involving shopping carts,
particularly heel and ankle impacts, are commonly associated with poor
braking response and momentum under heavy load \cite{smith2006,
vinck2019}. This highlights the need for intelligent control systems
capable of adjusting speed dynamically based on load conditions and
crowd density.

To address these challenges, this project proposes a Smart Autonomous
Shopping Cart Robot designed to aid users through adaptive navigation,
obstacle detection, human-aware motion control, and improved ergonomics.
The system aims to enhance safety, reduce physical strain, and improve
shopping efficiency in crowded retail environments.

Building on this foundation, this paper presents the simulation-based
implementation of the Smart Cart using ROS2 and Gazebo, validating
the proposed control architecture in a physics-accurate virtual
retail environment.

% ─── III. System Design ──────────────────────────────────────────────────────
\input{sections/system_dewsign_and_simulation}

% ─── IV. Non-Functional Requirements ─────────────────────────────────────────
\section{Non-Functional Requirements}

Non-functional requirements define the quality standards and operational
constraints the Smart Cart system must meet --- not what it does, but how
well it does it. These requirements directly influence the user experience,
safety, and deployment viability of the system.

\subsection{Performance Requirements}

The system must meet the following performance thresholds to function
safely and usefully in a real retail environment:

\begin{table}[!ht]
\caption{Non-Functional Requirements (Author developed)}
\label{tab:nfr}
\centering
\footnotesize
\begin{tabular}{p{2.0cm} p{1.8cm} p{2.8cm}}
\toprule
\textbf{Requirement} & \textbf{Target Value} & \textbf{Justification} \\
\midrule
Sensor-to-decision latency & $<$30\,ms & Cart must react to obstacles before reaching unsafe proximity at walking speed \\
Emergency stop response time & $<$200\,ms & Heel-impact prevention requires sub-200\,ms full brake from walking speed \\
UWB position update rate & $\geq$50\,Hz & Smooth follow-me requires frequent user position updates \\
Maximum following speed & 0.8\,m/s (unloaded) & Matches average comfortable shopping walking pace \\
Battery operational life & $\geq$3 hours continuous & Must cover a full extended shopping session without recharge \\
System boot-up time & $<$15 seconds & Acceptable wait time at cart collection point \\
Weight measurement accuracy & $\pm$200\,g & Sufficient for triggering load-based speed adjustments reliably \\
\bottomrule
\end{tabular}
\end{table}

\subsection{Safety Requirements}

Safety is the highest priority non-functional requirement. The cart
operates in a shared human environment and must never cause harm. The
following safety requirements are mandatory for deployment:
\begin{itemize}
  \item The cart shall never exceed 1.2\,m/s under any operating condition.
  \item The emergency stop mechanism shall operate via a hardware relay, independent of software --- cutting motor power even if the Raspberry Pi crashes.
  \item The cart shall detect and stop for obstacles lower than 20\,cm from the floor (child feet, bags) using IR sensors.
  \item The system shall prevent motion when basket load exceeds 40\,kg.
  \item If the UWB link to the user remote is lost for more than 3 seconds, the cart shall decelerate to a stop and wait.
  \item The cart shall not operate on slopes greater than 10° without triggering a low-speed safe mode.
  \item All exposed electrical components shall be rated IP42 minimum for indoor retail environments.
\end{itemize}

\subsection{Reliability and Availability}

The Smart Cart is a consumer-facing product deployed in a commercial
setting. Reliability failures have direct business and safety consequences.
\begin{itemize}
  \item Mean Time Between Failures (MTBF): The system shall operate for a minimum of 500 continuous shopping sessions before requiring scheduled maintenance.
  \item The software watchdog timer shall recover from non-critical ROS~2 node crashes within 2 seconds without user intervention.
  \item Battery management shall provide at least 60 seconds of safe operation warning before critical shutdown.
  \item All sensor failures shall be detected within one processing cycle (25\,ms) and trigger a safe-mode response.
\end{itemize}

\subsection{Usability Requirements}

The system must be accessible to all user classes without training,
including elderly shoppers and those unfamiliar with technology.
\begin{itemize}
  \item The entire user interface shall consist of two physical buttons only (Follow and Stop) --- no touchscreen, no app, no setup.
  \item LED status colors shall follow universally recognized conventions: green = safe/active, yellow = caution, red = stop/fault.
  \item The remote shall weigh no more than 80\,g and fit comfortably in one hand.
  \item A new user shall be able to begin using the cart within 30 seconds of receiving the remote, with no instruction beyond the LED color guide on the remote lanyard.
\end{itemize}

\subsection{Maintainability and Serviceability}

Mall operations staff must be able to maintain a fleet of Smart Carts
with minimal specialist expertise.
\begin{itemize}
  \item The battery pack shall be hot-swappable --- replaceable in under 3 minutes without tools.
  \item All sensor modules shall be individually replaceable without disassembling the main control enclosure.
  \item The system shall log all fault events with timestamps to an onboard SD card for post-session review by technical staff.
  \item Firmware updates shall be deployable over USB without requiring disassembly of the cart.
\end{itemize}

\subsection{Security Requirements}

Although Smart Cart does not handle personal data or payment information,
the remote-to-cart communication channel must be secured to prevent
malicious interference in a busy retail environment.
\begin{itemize}
  \item The UWB pairing between remote and cart shall use a unique session key generated at each power-on, preventing cross-cart interference or hijacking.
  \item No external network connectivity (WiFi, Bluetooth) shall be active during normal operation --- the system is fully self-contained.
  \item Firmware shall only accept signed update packages to prevent unauthorized code execution.
\end{itemize}


% ─── V. Experiment results ─────────────────────────────────────────────────
\section{Experimental Results}

\textit{This section is currently a draft. Simulation trials are 
ongoing and quantitative results will be finalized before submission.}

\subsection{Simulation Setup}
All experiments were conducted in Gazebo (version 11) with ROS2 
Humble. The simulated environment replicates a single retail aisle 
of 0.9\,m width and 10\,m length, with two static shelf obstacles 
and one dynamic pedestrian actor.

\subsection{Follow-Me Performance}
\textit{[DRAFT --- Data collection in progress.]}

Preliminary trials indicate the cart successfully maintains the 
target following distance of $0.8 \pm 0.2$\,m during straight-line 
tracking. Turn-following behavior is under evaluation.

\begin{table}[!ht]
\caption{Preliminary Follow-Me Results --- Gazebo Simulation 
         (Draft --- values to be confirmed)}
\label{tab:results_followme}
\centering
\footnotesize
\begin{tabular}{p{3.0cm} p{1.5cm} p{2.0cm}}
\toprule
\textbf{Test Scenario} & \textbf{Trials} & \textbf{Result} \\
\midrule
Straight-line follow (10\,m) & 5 & \textit{[pending]} \\
90° turn following           & 5 & \textit{[pending]} \\
Stop-and-resume behavior     & 5 & \textit{[pending]} \\
\bottomrule
\end{tabular}
\end{table}

\subsection{Obstacle Avoidance Performance}
\textit{[DRAFT --- Data collection in progress.]}

\begin{table}[!ht]
\caption{Preliminary Obstacle Avoidance Results --- Gazebo Simulation 
         (Draft --- values to be confirmed)}
\label{tab:results_obs}
\centering
\footnotesize
\begin{tabular}{p{3.0cm} p{1.5cm} p{2.0cm}}
\toprule
\textbf{Obstacle Type} & \textbf{Trials} & \textbf{Success Rate} \\
\midrule
Static shelf obstacle  & 10 & \textit{[pending]} \\
Dynamic pedestrian     & 10 & \textit{[pending]} \\
Narrow aisle (0.9\,m) & 5  & \textit{[pending]} \\
\bottomrule
\end{tabular}
\end{table}



% ─── VI. Feasibility Analysis ─────────────────────────────────────────────────
\subsection{Technical Feasibility}

All core technologies employed in the Smart Cart system are commercially
available, well-documented, and have been demonstrated in comparable
robotics research applications. The Raspberry Pi 4B running ROS~2 is a
proven platform for mobile robotics. UWB-based indoor positioning achieves
sub-30\,cm accuracy in retail-like environments, as documented in recent
literature. Omni-wheel holonomic platforms are a mature technology used
extensively in service robots. The primary technical risks are UWB
multipath interference in highly metallic shelf environments and real-time
control latency over the Arduino-to-RPi serial bridge --- both of which
have well-known mitigation strategies (antenna diversity and DMA serial
respectively).

\subsection{Economic Feasibility}

The following table provides a component-level bill of materials (BOM)
with estimated unit costs for a single prototype unit. Costs are based on
current retail pricing from distributor catalogues (RS Components, Mouser,
Amazon UK) as of 2025.

\begin{table}[!ht]
\caption{Bill of Materials --- Single Prototype Unit (Author developed)}
\label{tab:bom}
\centering
\footnotesize
\begin{tabular}{p{2.5cm} c r r p{1.3cm}}
\toprule
\textbf{Component} & \textbf{Qty} & \textbf{Unit (£)} & \textbf{Total (£)} & \textbf{Source} \\
\midrule
Raspberry Pi 4B (4GB) & 1 & 55 & 55 & RS Components \\
Arduino Mega 2560 & 1 & 38 & 38 & Arduino Official \\
DWM1000 UWB Module & 5 & 18 & 90 & Mouser \\
DC Motor + Omni Wheel Set & 4 & 22 & 88 & RobotShop \\
L298N Dual H-Bridge Driver & 2 & 8 & 16 & Amazon UK \\
HC-SR04 Ultrasonic Sensor & 2 & 4 & 8 & Amazon UK \\
Sharp IR Distance Sensor & 4 & 9 & 36 & RS Components \\
50kg Bar Load Cell + HX711 & 4 & 7 & 28 & Amazon UK \\
24V 20Ah LiPo Battery + BMS & 1 & 95 & 95 & Hobbyking \\
Buck Regulators (XL4016) & 2 & 6 & 12 & Amazon UK \\
RGB LED Strip (WS2812B) & 1 & 12 & 12 & Amazon UK \\
Remote PCB + Buttons + Case & 1 & 25 & 25 & Custom \\
Frame Mounts + Wiring + Misc & 1 & 40 & 40 & Various \\
\midrule
\textbf{TOTAL PROTOTYPE COST} & & & \textbf{£543} & \\
\bottomrule
\end{tabular}
\end{table}

For context, a standard commercial shopping cart costs approximately
£80--£150. The Smart Cart prototype at £543 represents a $\sim$4$\times$
cost multiplier for a single unit. At scale (500+ units), electronics
costs reduce significantly through bulk purchasing and custom PCB
fabrication, with an estimated production cost of £280--£320 per unit.
Given the operational benefits --- reduced staff assistance needs, reduced
cart damage, and improved shopper retention --- a payback period of 18--24
months is estimated for large-format retail deployments.

\subsection{Environmental Feasibility}

The system is designed for operation exclusively within controlled indoor
retail environments. Key environmental constraints and their mitigations
are summarized below:

\begin{table}[!ht]
\caption{Environmental Feasibility Assessment (Author developed)}
\label{tab:envfeas}
\centering
\footnotesize
\begin{tabular}{p{1.8cm} p{2.0cm} p{2.8cm}}
\toprule
\textbf{Environmental Factor} & \textbf{Constraint} & \textbf{Mitigation} \\
\midrule
Floor Surface & Smooth, level (max 3° ramp) & Omni-wheels handle minor slopes; ramp detection via IMU \\
Lighting & Moderate indoor ($>$200\,lux) & UWB and ultrasonic are light-independent \\
Crowd Density & Up to 2 people/m² & DWA planner with conservative speed limits \\
RF Interference & Dense WiFi / BLE environment & UWB at 3.5--6.5\,GHz; minimal overlap with WiFi \\
Temperature & 15--30°C standard retail & All components rated for 0--70°C operation \\
\bottomrule
\end{tabular}
\end{table}



% ─── VII. Testing Specifications ──────────────────────────────────────────────
\section{Testing Specifications}

Testing is structured in four progressive phases, from unit-level hardware
verification through to full-scenario integration tests in a simulated
retail environment. Each test specifies acceptance criteria that must be
passed before progressing to the next phase.

\subsection{Phase 1 --- Component Unit Tests}

\begin{table}[!ht]
\caption{Phase 1 --- Component Unit Tests (Author developed)}
\label{tab:phase1}
\centering
\footnotesize
\begin{tabular}{p{0.6cm} p{1.8cm} p{1.8cm} p{2.4cm}}
\toprule
\textbf{ID} & \textbf{Description} & \textbf{Method} & \textbf{Acceptance Criteria} \\
\midrule
T1.1 & UWB anchor ranging accuracy & Static tag at known positions (1\,m, 2\,m, 5\,m) & Error $<\pm$30\,cm at all ranges \\
T1.2 & Ultrasonic distance accuracy & Target board at 0.3, 0.5, 1.0, 2.0\,m & Error $<\pm$3\,cm; latency $<$25\,ms \\
T1.3 & IR sensor low-obstacle detection & Obstacle at 10, 40, 80\,cm floor level & Detected at all ranges within spec \\
T1.4 & Weight sensor calibration & Known masses: 0, 5, 10, 20, 40\,kg & Error $<\pm$200\,g across range \\
T1.5 & Motor direction and speed & PWM sweep 0--100\% all 4 motors & No stall; encoder counts consistent \\
T1.6 & Emergency stop response time & Remote X button press to motor stop & Full stop within 200\,ms \\
T1.7 & LiDAR 360° FoV Check & Place cart in clear area and observe ROS~2 \texttt{/scan} topic while rotating; verify basket frame occlusion & No blind spots or data gaps; static chassis frame must be filtered out \\
T1.8 & Camera Visual Odometry (VO) drift & Move cart using only camera with encoder disabled & Camera-based movement must have minimal drift ($<$10--15\,cm) over measured distance \\
\bottomrule
\end{tabular}
\end{table}

\subsection{Phase 2 --- Subsystem Integration Tests}

\begin{table}[!ht]
\caption{Phase 2 --- Subsystem Integration Tests (Author developed)}
\label{tab:phase2}
\centering
\footnotesize
\begin{tabular}{p{0.6cm} p{1.8cm} p{1.8cm} p{2.4cm}}
\toprule
\textbf{ID} & \textbf{Description} & \textbf{Method} & \textbf{Acceptance Criteria} \\
\midrule
T2.1 & ROS~2 topic publish rates & \texttt{ros2 topic hz} monitoring all nodes & All critical topics $>$40\,Hz \\
T2.2 & Serial bridge latency & Timestamp sensor data vs.\ ROS topic & End-to-end latency $<$30\,ms \\
T2.3 & EKF pose fusion stability & Walk user around cart perimeter 3$\times$ & No pose jumps $>$20\,cm; smooth trajectory \\
T2.4 & Obstacle zone classification & Approach cart with obstacle at varying speeds & Correct zone output at all thresholds \\
T2.5 & Speed adaptation to load & Test at 0, 10, 20, 40\,kg basket load & Speed reduces proportionally rules \\
T2.6 & Watchdog timer fault & Kill RPi process; observe Arduino & Motor cut within 500\,ms; red LED flash \\
\bottomrule
\end{tabular}
\end{table}

\subsection{Phase 3 --- Autonomous Navigation Tests}

\begin{table}[!ht]
\caption{Phase 3 --- Autonomous Navigation Tests (Author developed)}
\label{tab:phase3}
\centering
\footnotesize
\begin{tabular}{p{0.6cm} p{1.8cm} p{1.8cm} p{2.4cm}}
\toprule
\textbf{ID} & \textbf{Description} & \textbf{Method} & \textbf{Acceptance Criteria} \\
\midrule
T3.1 & Follow-me straight line & User walks 10\,m corridor; cart follows & Cart maintains 0.8\,m $\pm$ 0.2\,m following distance \\
T3.2 & Follow-me with turns & User makes 90° left and right turns & Cart completes turn within 1.5\,s with no collision \\
T3.3 & Obstacle avoidance --- static & Static box placed in cart path mid-follow & Cart re-routes without stopping $>$2\,s \\
T3.4 & Obstacle avoidance --- dynamic & Person walks across cart path at 1.0\,m/s & No contact; recovery to follow-path within 3\,s \\
T3.5 & Narrow aisle navigation & Aisle width of 0.9\,m (cart is 0.7\,m wide) & Cart passes without wall contact \\
\bottomrule
\end{tabular}
\end{table}

\subsection{Phase 4 --- Full Scenario / Performance Tests}

\begin{table}[!ht]
\caption{Phase 4 --- Full Scenario Performance Tests (Author developed)}
\label{tab:phase4}
\centering
\footnotesize
\begin{tabular}{p{0.6cm} p{1.8cm} p{1.8cm} p{2.4cm}}
\toprule
\textbf{ID} & \textbf{Description} & \textbf{Method} & \textbf{Acceptance Criteria} \\
\midrule
T4.1 & Full 2-hour battery endurance & Continuous operation, avg 0.5\,m/s, 15\,kg load & Battery $>$20\% at 2 hours; no fault interrupts \\
T4.2 & Simulated shopping route & User completes 20-min mall route with stops & $\geq$95\% successful follow-me engagement time \\
T4.3 & Crowded environment (10 people) & 10 actors walking in simulated mall space & Zero collisions; $\leq$3 emergency stops triggered \\
T4.4 & User stop-and-shop behaviour & User pauses for 30\,s, resumes walking & Cart holds position; resumes follow-me within 1\,s \\
T4.5 & Repeated E-Stop recovery & 10$\times$ stop/resume cycles in 5 minutes & All recoveries successful; no system fault logged \\
\bottomrule
\end{tabular}
\end{table}

% ─── VII. Risk Assessment ────────────────────────────────────────────────────
\section{Risk Assessment}

The following risk register identifies key technical, operational, and
safety risks associated with the Smart Cart system, along with their
likelihood, severity, and proposed mitigation strategies.

\begin{table}[!ht]
\caption{Risk Register --- Smart Cart System (Author developed)}
\label{tab:risk}
\centering
\footnotesize
\begin{tabular}{p{1.5cm} p{0.8cm} p{0.8cm} p{0.8cm} p{2.5cm}}
\toprule
\textbf{Risk} & \textbf{Like.} & \textbf{Sev.} & \textbf{Level} & \textbf{Mitigation} \\
\midrule
UWB multipath errors in metallic shelving & Med & High & HIGH & Deploy 4-anchor array; use EKF smoothing; test in representative environment \\
Cart collides with shopper leg or heel & Low & High & HIGH & Multi-zone obstacle detection; hardware E-stop relay; speed limits enforced in software \\
Battery depletion mid-session & Low & Med & MEDIUM & BMS low-battery alert at 20\%; geofenced return-to-base at 15\% \\
RPi thermal throttle under load & Med & Med & MEDIUM & Heatsink + 5\,V fan on RPi; monitor CPU temp via ROS~2 diagnostics node \\
Serial communication dropout between RPi and Arduino & Low & High & HIGH & Watchdog timer on Arduino; automatic re-connect in \texttt{sensor\_bridge} node \\
User walks faster than cart max speed & Med & Low & LOW & Cart switches to stationary follow mode; user notified via LED; max speed can be tuned \\
Children or pets triggering false obstacle stops & High & Low & MEDIUM & IR sensor height calibrated to ignore floor-level noise; adjustable stop thresholds \\
Cart remote lost or stolen & Med & Med & MEDIUM & Bluetooth pairing lock; remote-loss timeout (60\,s idle $\to$ auto park mode) \\
Omni-wheel performance on wet floors & Low & Med & MEDIUM & Speed reduction on incline (IMU); non-slip wheel treads; wet-floor warning in store policy \\
ROS~2 node crash causing uncontrolled motion & Low & High & HIGH & Watchdog timer; hardware relay E-stop; all motor commands time out if no update $>$300\,ms \\
\bottomrule
\end{tabular}
\end{table}

Risk levels are assessed using a standard 3$\times$3
likelihood-severity matrix. HIGH risks require mandatory mitigation prior
to deployment testing. MEDIUM risks require mitigation prior to full
public operation. LOW risks are monitored and reviewed at each testing
phase.


% ─── VIII. Future Improvements ───────────────────────────────────────────────
\section{Future Improvements}

The current Smart Cart specification represents a first-generation system
designed for practical deployment with commercially available components.
As the system matures through real-world testing and user feedback, several
enhancements are anticipated to improve capability, user experience, and
commercial viability.

\begin{table}[!ht]
\caption{Proposed Future Enhancements (Author developed)}
\label{tab:future}
\centering
\footnotesize
\begin{tabular}{p{1.8cm} p{2.5cm} p{2.5cm}}
\toprule
\textbf{Enhancement} & \textbf{Description} & \textbf{Expected Benefit} \\
\midrule
Shopping App Integration & Bluetooth link to shopper's phone to display shopping list and aisle navigation & Transforms cart into a personalised shopping assistant \\
Automatic Return-to-Base & When session ends and remote is docked, cart navigates autonomously back to charging station & Reduces staff workload; keeps carts charged and organised \\
Multi-Cart Fleet Management & Central dashboard for mall operators to monitor battery, location, and fault status of all carts & Improves operational efficiency at scale \\
Computer Vision Upgrade & Replace UWB tracking with camera-based human pose estimation (e.g.\ MediaPipe) & Removes need for user to carry remote; fully hands-free operation \\
Voice Feedback & Add a small speaker for audible alerts in noisy environments where LEDs may be missed & Improves accessibility for visually impaired shoppers \\
\bottomrule
\end{tabular}
\end{table}

Of these enhancements, the Computer Vision Upgrade and Automatic
Return-to-Base represent the highest-priority upgrades for a
second-generation prototype, as they directly address the primary
technical limitation of the current design --- its reliance on
short-range ultrasonic sensing for obstacle detection.


% ─── Appendix D: Use of Generative AI ────────────────────────────────────────
\appendix
\section{Use of Generative AI}
\textit{This appendix has been updated to reflect AI tool usage 
during CW2 simulation implementation and report drafting.}

In accordance with Middlesex University's assessment guidelines and the
module requirement to declare AI usage, this appendix discloses how and
where generative AI tools were used in the preparation of this document.

\subsection{How AI Was Used}

The team used Claude (Anthropic) as a writing and structuring aid
throughout the preparation of this report. AI assistance was used in the
following ways:

\begin{table}[!ht]
\caption{AI Usage Disclosure by Section (Author developed)}
\label{tab:ai}
\centering
\footnotesize
\begin{tabular}{p{1.8cm} p{2.2cm} p{2.5cm}}
\toprule
\textbf{Section} & \textbf{AI Role} & \textbf{Human Contribution} \\
\midrule
Introduction and Problem Statement & Draft structure and academic language polishing & Original problem observations, research citations, and framing \\
Hardware Specifications (Section~III) & Generating specification tables with realistic component values & Team verified all components against real datasheets and design diagram \\
Software Architecture (Section~III) & Drafting ROS~2 node graph and communication protocol tables & Team defined the actual node structure, sensor choices, and control logic \\
Feasibility and Cost Analysis (Section~V) & Generating BOM table with estimated costs & Team reviewed all prices against current supplier listings and adjusted \\
Testing Specifications (Section~VI) & Drafting test case structure and acceptance criteria & Team defined test scenarios based on actual design goals \\
Risk Assessment (Section~VII) & Populating risk register template & Team identified risks based on real design knowledge and class content \\
\bottomrule
\end{tabular}
\end{table}

\subsection{Human Oversight and Verification}

All AI-generated content was reviewed, edited, and verified by team
members before inclusion in this report. The core design decisions ---
including the choice of UWB positioning over camera-based tracking, the
selection of omni-wheels, the Raspberry Pi + Arduino dual-processor
architecture, and the remote-control interface --- were developed
independently by the team through group discussion, class knowledge, and
the design sketch produced during the ideation phase. AI was used as a
productivity and language tool, not as a design decision-maker.

Team members take full academic responsibility for the content of this
document.

\subsection{AI Tools Used}
\begin{itemize}
  \item Claude (Anthropic) --- document drafting, structuring, and table generation
  \item No AI-based image generation tools were used; all diagrams are author-developed
\end{itemize}

% ─── Appendix B: Design Thinking & Storyboard ────────────────────────────────
\section{Design Thinking \& Storyboard}

This section documents the design thinking process used by the team to
arrive at the final Smart Cart concept. Design thinking was applied across
five stages: Empathize, Define, Ideate, Prototype, and Test. Each stage
shaped key decisions about the system's hardware configuration, user
interaction model, and safety approach.

\subsection{Empathize --- Understanding the Problem}

The team began by observing shopping behavior in a large retail mall
environment. Key pain points identified through observation and informal
user interviews included:
\begin{itemize}
  \item Elderly shoppers struggled to push heavily loaded carts, especially on slight inclines near entrances.
  \item Parents with young children frequently lose control of carts when distracted.
  \item Narrow aisles in peak hours caused collisions between carts and other shoppers, particularly at ankle and heel height.
  \item Cart return was often neglected, causing trolleys to roll freely in car parks.
\end{itemize}

These direct observations motivated the core design goal: a cart that
removes physical burden from the user while actively preventing collisions
in dynamic environments.

\subsection{Define --- Problem Statement}

Based on the empathy stage, the team defined the central problem as:

\begin{quote}
\textit{``Shoppers in crowded retail environments need a smarter cart
that follows them autonomously, detects obstacles, and adapts to load ---
so they can shop safely without physical strain.''}
\end{quote}

\subsection{Ideate --- Design Concepts Considered}

The team generated three competing design concepts before converging on
the final solution:

\begin{table}[!ht]
\caption{Design Concepts Considered During Ideation (Author developed)}
\label{tab:ideate}
\centering
\footnotesize
\begin{tabular}{p{1.2cm} p{2.3cm} p{3.0cm}}
\toprule
\textbf{Concept} & \textbf{Description} & \textbf{Why Rejected / Adopted} \\
\midrule
Concept A & GPS-guided outdoor cart with motorized wheels & GPS unreliable indoors; rejected for indoor-only scope \\
Concept B & Camera-only follow system (no UWB) & Camera occlusion in crowds caused unreliable tracking; rejected \\
Concept C (Chosen) & UWB-based follow-me with omni-wheels, ultrasonic + IR sensors & Most reliable indoors; handles occlusion; low cost; ADOPTED \\
\bottomrule
\end{tabular}
\end{table}

% ─── References ──────────────────────────────────────────────────────────────
\bibliographystyle{IEEEtran}
\bibliography{references}
\end{document}
